\chapter{Regime Modeling}
\label{ch:regimes}

\begin{projectconnection}
In Phase~4A (Task~17), you fit a Gaussian Mixture Model on macro features to
classify market regimes and decompose your signal's performance by regime.
This chapter teaches GMM regime classification, the Two Sigma template for
interpreting regimes, ADWIN for online drift detection, Markov-switching
models for regime persistence, and---critically---the Asness et al.\ (2017)
caveat on why regime-conditional strategies often disappoint in practice.
\end{projectconnection}

% ══════════════════════════════════════════════════════════════
\section{Why Regime-Awareness Matters}
\label{sec:regimes:why}
% ══════════════════════════════════════════════════════════════

A signal that earns $\IC = 0.04$ unconditionally might earn $\IC = 0.10$ in
calm markets and $\IC = -0.03$ during crises.  If the positive performance
comes entirely from one market regime, the signal is \emph{fragile}: the
next crisis could wipe out years of steady gains.  Regime decomposition
answers a simple but important question: \emph{is your signal's Sharpe
driven by one regime, or is it broad-based?}

Three practical uses of regime modeling in your project:

\begin{enumerate}[leftmargin=1.5em]
  \item \textbf{Diagnostic.}  Decompose backtest Sharpe by regime.  A signal
    that works only in steady-state periods is less valuable than one that
    works across regimes.
  \item \textbf{Risk management.}  Scale position sizes down when the model
    assigns high probability to a crisis regime.
  \item \textbf{Feature enrichment.}  Feed regime probabilities as
    additional features into your LightGBM model (Chapter~\ref{ch:trees}).
\end{enumerate}

\begin{intuition}[Regimes as Weather]
Think of regimes like weather seasons.  A farmer who only knows the
\emph{average} annual rainfall will plant the same crop year-round.  One who
distinguishes wet and dry seasons adapts.  Your signal is the crop;
regimes are the seasons.  The goal is not to predict regimes (that is
hard)---it is to \emph{check whether your crop survives all seasons}.
\end{intuition}

% ══════════════════════════════════════════════════════════════
\section{Gaussian Mixture Models for Regime Classification}
\label{sec:regimes:gmm}
% ══════════════════════════════════════════════════════════════

\begin{prereq}[Multivariate Gaussian and Mixture Models]
A $d$-dimensional Gaussian density is
\[
  \N(\bx \mid \bm{\mu}, \bm{\Sigma})
  = \frac{1}{(2\pi)^{d/2} |\bm{\Sigma}|^{1/2}}
    \exp\!\Bigl(-\tfrac{1}{2}(\bx - \bm{\mu})^\top \bm{\Sigma}^{-1}
    (\bx - \bm{\mu})\Bigr),
\]
where $\bm{\mu} \in \R^d$ is the mean and
$\bm{\Sigma} \in \R^{d \times d}$ is the covariance matrix
(symmetric positive-definite).  A \emph{mixture model} says the data
come from several such Gaussians, each chosen with some probability.
If this is unfamiliar, review Bishop (2006), Ch.~2 and~9.
\end{prereq}

\begin{definition}[Gaussian Mixture Model]
A $K$-component Gaussian Mixture Model (GMM) models the density of
observations $\bx \in \R^d$ as
\begin{equation}
  p(\bx) = \sum_{k=1}^{K} \pi_k \; \N(\bx \mid \bm{\mu}_k, \bm{\Sigma}_k),
  \label{eq:gmm}
\end{equation}
where:
\begin{itemize}[leftmargin=1.5em]
  \item $K$ is the number of mixture components (regimes),
  \item $\pi_k \geq 0$ is the \textbf{mixing weight} for component $k$,
        with $\sum_{k=1}^K \pi_k = 1$,
  \item $\bm{\mu}_k \in \R^d$ is the \textbf{mean} of component $k$,
  \item $\bm{\Sigma}_k \in \R^{d \times d}$ is the \textbf{covariance}
        of component $k$.
\end{itemize}
\end{definition}

The parameters $\btheta = \{\pi_k, \bm{\mu}_k, \bm{\Sigma}_k\}_{k=1}^K$
are estimated by maximum likelihood.  Because the log-likelihood
\[
  \ell(\btheta) = \sum_{t=1}^{T} \log \sum_{k=1}^{K} \pi_k \;
  \N(\bx_t \mid \bm{\mu}_k, \bm{\Sigma}_k)
\]
has no closed-form solution (the log wraps a sum), we use the
Expectation--Maximization (EM) algorithm
(Dempster, Laird, and Rubin, 1977).

% ── EM Algorithm ──
\subsection{The EM Algorithm}
\label{sec:regimes:em}

EM alternates between two steps until convergence:

\begin{keyidea}[EM for Gaussian Mixtures]
\textbf{E-step} (soft labels).  For each observation $\bx_t$ and component
$k$, compute the \emph{responsibility}:
\begin{equation}
  \gamma_{tk} = \frac{\pi_k \; \N(\bx_t \mid \bm{\mu}_k, \bm{\Sigma}_k)}
    {\sum_{j=1}^{K} \pi_j \; \N(\bx_t \mid \bm{\mu}_j, \bm{\Sigma}_j)}.
  \label{eq:estep}
\end{equation}
$\gamma_{tk}$ is the posterior probability that observation $t$ belongs to
component $k$.

\textbf{M-step} (update parameters).  Using the responsibilities as weights:
\begin{align}
  N_k &= \sum_{t=1}^{T} \gamma_{tk}, \label{eq:mstep-n}\\[4pt]
  \bm{\mu}_k^{\text{new}} &= \frac{1}{N_k} \sum_{t=1}^{T} \gamma_{tk}\,
    \bx_t, \label{eq:mstep-mu}\\[4pt]
  \bm{\Sigma}_k^{\text{new}} &= \frac{1}{N_k} \sum_{t=1}^{T}
    \gamma_{tk}\,(\bx_t - \bm{\mu}_k^{\text{new}})
    (\bx_t - \bm{\mu}_k^{\text{new}})^\top, \label{eq:mstep-sigma}\\[4pt]
  \pi_k^{\text{new}} &= \frac{N_k}{T}. \label{eq:mstep-pi}
\end{align}
$N_k$ is the effective number of points assigned to component $k$.

Repeat E-step and M-step until $|\ell^{(i+1)} - \ell^{(i)}| < \epsilon$.
EM guarantees $\ell(\btheta)$ is non-decreasing at every iteration
(Dempster, Laird, and Rubin, 1977).
\end{keyidea}

% ── Model selection ──
\subsection{Choosing $K$: BIC}
\label{sec:regimes:bic}

More components always improve in-sample fit, so you cannot pick $K$ by
maximizing the log-likelihood.  The Bayesian Information Criterion penalizes
complexity:
\begin{equation}
  \text{BIC} = -2\,\ell(\hat{\btheta}) + p \log T,
  \label{eq:bic}
\end{equation}
where $p$ is the number of free parameters and $T$ is the sample size.
For a $K$-component GMM in $d$ dimensions with full covariance matrices,
$p = K(1 + d + d(d+1)/2) - 1$.  Choose the $K$ that
\emph{minimizes} BIC.  Two Sigma used cross-validated log-likelihood with
similar results.  For your project, $K \in \{2, 3, 4\}$; $K = 3$ or $4$
is typical.

% ── Feature choice ──
\subsection{Feature Selection for Regime GMM}
\label{sec:regimes:features}

The regime GMM takes macro observables, not your alpha features.  Typical
choices:

\begin{center}
\begin{tabular}{lll}
  \toprule
  Feature & Source & Captures \\
  \midrule
  VIX level & CBOE & Equity risk appetite \\
  IG credit spread & ICE BofA & Credit stress \\
  10y--2y term slope & Treasury & Growth/policy expectations \\
  DXY (USD index) & Fed & Dollar regime \\
  Realized cross-asset corr. & Computed & Contagion / diversification \\
  \bottomrule
\end{tabular}
\end{center}

Standardize all features to zero mean and unit variance before fitting.
This prevents the GMM from being dominated by whichever feature has the
largest scale.

\begin{warning}[Do Not Use Your Alpha Features as Regime Inputs]
If you feed your risk-cube signals into the regime GMM \emph{and then}
condition on regimes, you introduce circularity.  The regime model should
use independent macro variables.  Your alpha features go into the
LightGBM prediction model; the regime model provides an orthogonal lens
for decomposing performance.
\end{warning}

% ── Worked example ──
\begin{workedexample}{One EM Iteration on a 2D, 3-Component GMM}
Suppose $d = 2$ features (VIX, credit spread), $K = 3$ components, and
$T = 5$ observations.  After initialization, the current parameters are:

\[
  \pi_1 = 0.5, \quad \pi_2 = 0.3, \quad \pi_3 = 0.2.
\]
\[
  \bm{\mu}_1 = \begin{pmatrix} 14 \\ 1.0 \end{pmatrix}, \quad
  \bm{\mu}_2 = \begin{pmatrix} 22 \\ 2.5 \end{pmatrix}, \quad
  \bm{\mu}_3 = \begin{pmatrix} 35 \\ 5.0 \end{pmatrix}.
\]
All covariances initialized to $\bm{\Sigma}_k = \mathbf{I}_2$ for
simplicity.

\textbf{E-step for $\bx_1 = (15, 1.2)^\top$:}

Compute unnormalized responsibilities:
\begin{align*}
  \tilde{\gamma}_{1,1} &= \pi_1 \cdot \N(\bx_1 \mid \bm{\mu}_1,
    \mathbf{I}) = 0.5 \times \N\bigl((15,1.2)^\top \mid (14,1.0)^\top,
    \mathbf{I}\bigr). \\
  \|\bx_1 - \bm{\mu}_1\|^2 &= (15-14)^2 + (1.2-1.0)^2 = 1.04. \\
  \N(\cdot) &= \frac{1}{2\pi} \exp(-1.04/2) = 0.0942. \\
  \tilde{\gamma}_{1,1} &= 0.5 \times 0.0942 = 0.0471.
\end{align*}
Similarly, $\|\bx_1 - \bm{\mu}_2\|^2 = (15-22)^2 + (1.2-2.5)^2 = 50.69$,
giving $\tilde{\gamma}_{1,2} = 0.3 \times 1.32 \times 10^{-12} \approx 0$.
And $\|\bx_1 - \bm{\mu}_3\|^2 = 414.44$, giving
$\tilde{\gamma}_{1,3} \approx 0$.

\textbf{Normalize:}
\[
  \gamma_{1,1} = \frac{0.0471}{0.0471 + 0 + 0} \approx 1.00, \quad
  \gamma_{1,2} \approx 0, \quad \gamma_{1,3} \approx 0.
\]

Observation $\bx_1$ is assigned almost entirely to component~1 (low-VIX,
low-spread = steady state).  Repeat for all $T$ observations to complete
the E-step, then update parameters via Eqs.~\eqref{eq:mstep-n}--\eqref{eq:mstep-pi}.

\textbf{Takeaway.}  Points far from a component's mean get
near-zero responsibility.  This is how the GMM ``self-organizes'' into
clusters without hard labels.
\end{workedexample}

% ══════════════════════════════════════════════════════════════
\section{The Two Sigma Regime Template}
\label{sec:regimes:twosigma}
% ══════════════════════════════════════════════════════════════

Two Sigma applied a GMM to 17 factor returns and identified four market
conditions that recur historically.  Their template provides a useful
vocabulary for interpreting your own GMM output.

\begin{keyidea}[Four Macro Regimes (Two Sigma, 2021)]
\begin{center}
\begin{tabular}{p{2.8cm} p{5.5cm} p{5.2cm}}
  \toprule
  \textbf{Regime} & \textbf{Characteristics} & \textbf{Asset behavior} \\
  \midrule
  \textbf{Steady State} &
    Low VIX, tight spreads, moderate growth.  Dominates the 2010s decade. &
    Equity, credit, and most style factors earn positive returns. \\[4pt]
  \textbf{Crisis} &
    VIX spikes, spreads widen, cross-asset correlations surge.
    1987, 2008~GFC, 2020~COVID. &
    Equity and credit suffer large drawdowns; defensive factors outperform. \\[4pt]
  \textbf{Inflation} &
    Rising CPI, steepening yield curve, commodity strength.
    Concentrated in 1970s--early 1980s. &
    Local inflation factor posts double-digit mean returns; equities and
    bonds underperform. \\[4pt]
  \textbf{Walking on Ice} &
    Equity returns above average but volatility $\sim$1.6~pp above
    long-run mean.  Often precedes or follows crises. &
    Positive equity returns mask elevated fragility; bubble risk. \\
  \bottomrule
\end{tabular}
\end{center}
\end{keyidea}

\subsection{Regime-Conditional Sharpe Decomposition}
\label{sec:regimes:conditional-sharpe}

Once you assign regime probabilities to each date in your backtest, compute
regime-conditional Sharpe ratios:
\begin{equation}
  \SR_k = \frac{\bar{r}_k}{\sigma_k}, \quad
  \text{where } \bar{r}_k = \frac{\sum_t \gamma_{tk}\, r_t}
  {\sum_t \gamma_{tk}}, \quad
  \sigma_k^2 = \frac{\sum_t \gamma_{tk}\,(r_t - \bar{r}_k)^2}
  {\sum_t \gamma_{tk}}.
  \label{eq:regime-sharpe}
\end{equation}
Report these alongside your unconditional Sharpe.  The presentation
(Chapter~\ref{ch:presenting}) should include a table like:

\begin{center}
\begin{tabular}{lcccc}
  \toprule
  & Steady State & Crisis & Inflation & Walking on Ice \\
  \midrule
  $\SR$ & 1.20 & $-0.30$ & 0.45 & 0.85 \\
  Fraction of days & 55\% & 10\% & 15\% & 20\% \\
  \bottomrule
\end{tabular}
\end{center}

If $\SR$ is positive only in Steady State (55\% of days), your signal is
fragile.  If positive across three or four regimes, it is robust.

\begin{prereq}[Sharpe Ratio and Annualization]
The Sharpe ratio $\SR = \bar{r}/\sigma$ measures risk-adjusted return.
For daily data, annualize via $\SR_{\text{annual}} = \SR_{\text{daily}}
\times \sqrt{252}$.  This chapter computes regime-conditional Sharpe
ratios on daily strategy returns, then annualizes.  See
Chapter~\ref{ch:backtesting} for the full backtesting framework.
\end{prereq}

% ══════════════════════════════════════════════════════════════
\section{ADWIN for Concept Drift Detection}
\label{sec:regimes:adwin}
% ══════════════════════════════════════════════════════════════

The GMM approach is \emph{offline}: you fit the model to all available data
and then label each date.  In production, you need \emph{online} detection
of distribution shifts.  ADWIN (Adaptive Windowing; Bifet and Gavald\`{a},
2007) provides exactly this.

\begin{definition}[ADWIN Algorithm]
ADWIN maintains a variable-length window $W$ of recent observations.  At
each time step, it tests all possible partitions of $W$ into two contiguous
sub-windows $W_0$ (older) and $W_1$ (newer).  A drift is declared when the
absolute difference between the sub-window means exceeds a threshold
derived from the Hoeffding bound:
\begin{equation}
  |\hat{\mu}_{W_0} - \hat{\mu}_{W_1}| \geq \epsilon_{\text{cut}},
  \quad \text{where } \epsilon_{\text{cut}}
  = \sqrt{\frac{1}{2m}\cdot \ln \frac{4|W|}{\delta}},
  \label{eq:adwin}
\end{equation}
$m = (1/|W_0| + 1/|W_1|)^{-1}$ is the harmonic mean of the sub-window
sizes, $|W|$ is the total window length, and $\delta \in (0,1)$ is a
confidence parameter controlling sensitivity.  When drift is detected,
all observations in $W_0$ are discarded.
\end{definition}

\textbf{How to use ADWIN in your project:}
\begin{enumerate}[leftmargin=1.5em]
  \item Feed your model's \emph{prediction errors}
        ($\hat{y}_t - y_t$) into ADWIN.
  \item When ADWIN fires, it means the error distribution has shifted---your
        model's relationship with the data has changed.
  \item Trigger a model retrain on the post-drift window only.
\end{enumerate}

The parameter $\delta$ controls the trade-off: $\delta = 0.01$ (more
sensitive, more false positives, faster detection) vs.\ $\delta = 0.0001$
(fewer false positives, slower detection).  For daily financial data with
$\sim$252 observations per year, $\delta = 0.002$ is a reasonable starting
point.

\begin{intuition}[ADWIN as a Smoke Detector]
ADWIN is like a smoke detector for your model.  It does not tell you
\emph{what} regime you are in (the GMM does that).  It tells you that
something has \emph{changed}---your model's errors suddenly look different
from what they were.  When the alarm sounds, you investigate: has the
market entered a new regime, or did your feature pipeline break?
\end{intuition}

% ══════════════════════════════════════════════════════════════
\section{Markov-Switching Models}
\label{sec:regimes:markov}
% ══════════════════════════════════════════════════════════════

\begin{prereq}[Hidden Markov Models]
A Hidden Markov Model (HMM) has two layers: (1)~a latent state sequence
$s_t \in \{1, \ldots, K\}$ evolving as a Markov chain, and (2)~observed
emissions $\bx_t$ whose distribution depends on $s_t$.  The key
assumption: $s_t$ depends on $s_{t-1}$ only (Markov property).  If
unfamiliar, review Murphy (2012), Ch.~17.
\end{prereq}

Hamilton (1989) introduced Markov-switching regression to economics,
modeling U.S.\ GNP growth as an autoregression whose mean shifts between
two states: expansion and contraction.  The model is:
\begin{equation}
  y_t = \mu_{s_t} + \sum_{i=1}^{p} \phi_i (y_{t-i} - \mu_{s_{t-i}})
  + \varepsilon_t, \quad \varepsilon_t \sim \N(0, \sigma^2),
  \label{eq:hamilton-ms}
\end{equation}
where $s_t \in \{0, 1\}$ is the latent state with transition matrix:
\begin{equation}
  \mathbf{P} = \begin{pmatrix}
    p_{00} & p_{01} \\
    p_{10} & p_{11}
  \end{pmatrix}, \quad
  p_{ij} = \Pr(s_t = j \mid s_{t-1} = i), \quad
  p_{i0} + p_{i1} = 1.
  \label{eq:transition}
\end{equation}

\begin{keyresult}[Hamilton (1989): Markov-Switching GNP]
Applied to U.S.\ quarterly GNP growth (1951Q1--1984Q4) with an AR(4)
specification, Hamilton estimated:
\begin{itemize}[leftmargin=1.5em]
  \item Mean growth in expansion ($s_t = 1$): $\hat{\mu}_1 = 1.16\%$
        per quarter.
  \item Mean growth in recession ($s_t = 0$): $\hat{\mu}_0 = -0.36\%$
        per quarter.
  \item Persistence: $\hat{p}_{11} = 0.90$ (expansions last on average
        $1/(1 - 0.90) \approx 10$ quarters),
        $\hat{p}_{00} = 0.75$ (recessions last $\approx$4 quarters).
\end{itemize}
The model's filtered probabilities closely matched NBER recession
dates without using them as inputs.
\end{keyresult}

\subsection{GMM vs.\ Markov-Switching: When to Use Which}
\label{sec:regimes:comparison}

\begin{center}
\begin{tabular}{p{3.5cm} p{5cm} p{5cm}}
  \toprule
  & \textbf{GMM} & \textbf{Markov-Switching} \\
  \midrule
  Temporal structure &
    None.  Each $\bx_t$ classified independently. &
    Markov chain: $s_t$ depends on $s_{t-1}$. \\[4pt]
  Regime persistence &
    Not modeled.  Can oscillate between regimes at high frequency. &
    Explicitly modeled via transition probabilities. \\[4pt]
  Estimation &
    EM on mixture (fast, scales well). &
    Hamilton filter (slower, requires forward--backward recursion). \\[4pt]
  Output &
    Posterior probabilities $\gamma_{tk}$ per observation. &
    Filtered $\Pr(s_t = k \mid \bx_{1:t})$ and smoothed
    $\Pr(s_t = k \mid \bx_{1:T})$. \\[4pt]
  Python &
    \texttt{sklearn.mixture.GaussianMixture} &
    \texttt{statsmodels.tsa.regime\_switching} \\[4pt]
  Best for &
    Quick regime labeling; diagnostic decomposition. &
    When regime persistence matters; macro modeling. \\
  \bottomrule
\end{tabular}
\end{center}

\begin{keyidea}[Use GMM for Your Project]
GMM treats each observation independently; Markov-switching models regime
persistence.  For your project, you are performing a \emph{post-hoc
diagnostic}: ``did my signal's Sharpe come from one regime or many?''
You are not forecasting regimes in real time.  For this purpose, GMM is
simpler, faster, and sufficient.  If you later want to build a regime
forecasting overlay for position sizing, switch to Markov-switching.
\end{keyidea}

% ══════════════════════════════════════════════════════════════
\section{The Factor Timing Caveat}
\label{sec:regimes:timing-caveat}
% ══════════════════════════════════════════════════════════════

Before you get excited about regime-conditional strategies, read the
cautionary evidence.

\begin{keyresult}[Asness, Chandra, Ilmanen, and Israel (2017)]
In ``Contrarian Factor Timing is Deceptively Difficult,'' the authors test
whether buying factors after drawdowns (contrarian timing via value spreads)
improves multi-factor portfolios.  Key findings:
\begin{itemize}[leftmargin=1.5em]
  \item For a 3-factor portfolio (HML + UMD + BAB), value-spread timing
        \emph{decreases} the Sharpe ratio from 0.55 to 0.54.
  \item For individual factors, timing helps momentum (Sharpe 0.26
        $\to$ 0.33) but hurts defensive/low-beta (Sharpe 0.53 $\to$ 0.40).
  \item The timing signal is correlated with the value factor itself, so
        in a portfolio that already holds value, timing adds redundant and
        intermittent exposure rather than genuine new information.
  \item Results are reported gross of transaction costs.  Net of costs,
        the case for timing weakens further.
\end{itemize}
The authors conclude that contrarian factor timing is ``a weak addition for
long-term investors holding well-diversified factors including value.''
\end{keyresult}

\subsection{Implications for Your Project}
\label{sec:regimes:implications}

The Asness et al.\ result does \emph{not} say regime modeling is useless.
It says:

\begin{enumerate}[leftmargin=1.5em]
  \item \textbf{Conditional $>$ unconditional is the wrong comparison.}
    You must compare the regime-conditional strategy to the unconditional
    strategy \emph{after accounting for any additional factor exposures} the
    conditioning introduces.  If your regime overlay tilts toward value or
    momentum, you are adding a factor bet, not alpha.
  \item \textbf{Document honestly.}  Report both conditional and
    unconditional Sharpe.  If the conditional strategy wins, show that it
    is not explained by loading on known factors.
  \item \textbf{The bar is high.}  Even sophisticated practitioners at AQR
    found that timing adds little.  If your regime decomposition shows that
    the signal works only in one regime, the honest conclusion may be to
    report that fragility rather than to build a timing strategy around it.
\end{enumerate}

\begin{warning}[Regimes Are Only Known in Hindsight]
Regime-conditional strategies look better in backtests than in practice
because the regime is only known with certainty in hindsight.  Your GMM
assigns \emph{probabilities}, not hard labels.  Today's observation might
be 60\% Steady State and 40\% Walking on Ice.  If you build a strategy
that goes fully risk-on in Steady State and risk-off otherwise, you are
implicitly assuming perfect regime classification.  In practice: (a)~use
soft probabilities, not hard assignments; (b)~test with lagged regime
labels (yesterday's classification applied today); (c)~measure how
often the GMM changes its mind on the most recent 20 days of data when
one new observation arrives.
\end{warning}

% ══════════════════════════════════════════════════════════════
\section{Practical Implementation Checklist}
\label{sec:regimes:checklist}
% ══════════════════════════════════════════════════════════════

For Task~17 in your project plan, follow this sequence:

\begin{enumerate}[leftmargin=1.5em]
  \item \textbf{Assemble macro features.}  Pull VIX, IG credit spread,
    10y--2y slope, DXY, and realized 60-day cross-asset correlation.
    Standardize to $z$-scores.
  \item \textbf{Fit GMMs for $K \in \{2, 3, 4, 5\}$.}  Use
    \texttt{sklearn.mixture.GaussianMixture} with
    \texttt{covariance\_type='full'} and \texttt{n\_init=10} (multiple
    random restarts to avoid local optima).  Record BIC for each $K$.
  \item \textbf{Select $K$.}  Pick the $K$ that minimizes BIC.  Inspect
    the regimes manually: do they correspond to recognizable market
    conditions?  If $K = 3$ and $K = 4$ have similar BIC, prefer the one
    with more interpretable clusters.
  \item \textbf{Label each date.}  Assign soft probabilities
    $\gamma_{tk}$, not hard labels.  Store as columns in your feature
    DataFrame.
  \item \textbf{Decompose signal performance.}  Compute
    regime-conditional Sharpe (Eq.~\ref{eq:regime-sharpe}), IC, and
    hit rate.  Report in a table.
  \item \textbf{Test regime overlay.}  If desired, scale positions by
    $\gamma_{t,\text{steady}}$ (probability of being in the favorable
    regime).  Compare Sharpe to the unconditional strategy.  Apply the
    Asness et al.\ sanity check: does the improvement survive controlling
    for additional factor exposures?
  \item \textbf{Set up ADWIN.}  Feed prediction residuals into ADWIN
    with $\delta = 0.002$.  Log all detected drift points.  Check whether
    they align with GMM regime transitions.
\end{enumerate}

% ══════════════════════════════════════════════════════════════
\section{Summary}
\label{sec:regimes:summary}
% ══════════════════════════════════════════════════════════════

\begin{center}
\begin{tabular}{p{4cm} p{9.5cm}}
  \toprule
  \textbf{Concept} & \textbf{Key Takeaway} \\
  \midrule
  GMM regime classification &
    $p(\bx) = \sum_k \pi_k \N(\bx \mid \bm{\mu}_k, \bm{\Sigma}_k)$.
    EM alternates E-step (soft labels) and M-step (update parameters).
    Select $K$ by BIC. \\[4pt]
  Two Sigma template &
    Four regimes---Steady State, Crisis, Inflation, Walking on Ice---fit
    to factor returns via GMM. \\[4pt]
  ADWIN &
    Online drift detection via adaptive windowing and Hoeffding bounds.
    Apply to prediction errors to trigger retraining. \\[4pt]
  Markov-switching &
    Hamilton (1989): HMM with regime-dependent parameters.  Models
    persistence.  Use GMM for post-hoc decomposition; Markov-switching
    for real-time regime forecasting. \\[4pt]
  Factor timing caveat &
    Asness et al.\ (2017): contrarian timing of factors via value spreads
    barely helps (3-factor Sharpe 0.55 $\to$ 0.54) and can hurt
    diversification.  Regime conditioning is diagnostic first, strategy
    second. \\[4pt]
  Your project &
    Phase~4A, Task~17: fit GMM on macro features, decompose signal by
    regime, report conditional Sharpe honestly.  Use soft probabilities,
    test with lagged labels, control for factor exposure overlap. \\
  \bottomrule
\end{tabular}
\end{center}
